\documentclass[11pt,a4paper]{article}

\usepackage[utf8]{inputenc}
\usepackage[T1]{fontenc}
\usepackage[ngerman]{babel}
\usepackage{lmodern}
\usepackage{geometry}
\usepackage{xcolor}
\usepackage{enumitem}
\usepackage{hyperref}

\geometry{margin=2.5cm}
\setlength{\parindent}{0pt}
\setlength{\parskip}{0.75em}
\hypersetup{
  colorlinks=true,
  linkcolor=blue,
  urlcolor=blue
}

\newcommand{\folie}[2]{%
  \section*{Folie #1: #2}
}

\title{Vortragstext zur Präsentation\\
\large Teddybär als Baby-Weinen-Alarm}
\author{Projektvorstellung}
\date{\today}

\begin{document}

\maketitle

\section*{Hinweis}
Dieser Text ist als Sprechtext gedacht. Er muss nicht wortwörtlich
abgelesen werden, sondern kann als Orientierung für die Präsentation
verwendet werden.

\folie{1}{Titelfolie}
Guten Tag, ich stelle heute meinen Prototypen \glqq Teddybär als
Baby-Weinen-Alarm\grqq{} vor.

In dem Projekt geht es um ein interaktives Textil. Ein normales,
vertrautes Objekt, nämlich ein Kuschelbär, wird mit Sensorik und
Software erweitert. Der Bär soll laute Geräusche oder mögliches
Babyweinen erkennen und diese Information an eine App weitergeben.

Technisch besteht der Prototyp aus einem ESP32 mit Mikrofon, einem
Go-Server zur Auswertung der Messwerte und einer Flutter-App, die den
aktuellen Zustand für die Eltern sichtbar macht.

\folie{2}{Ausgangsidee}
Die Ausgangsidee war, Technik möglichst unauffällig in ein vertrautes
Objekt zu integrieren.

Ein Teddybär ist weich, bekannt und wirkt nicht wie ein technisches
Gerät. Genau das ist bei interaktiven Textilien interessant: Die
Elektronik verschwindet im Objekt, aber das Objekt bekommt trotzdem eine
digitale Funktion.

Der Bär nimmt Geräusche in seiner Umgebung wahr. Die Messwerte werden an
einen Server geschickt. Dort wird entschieden, ob alles ruhig ist, ob es
sehr laut ist oder ob mögliches Babyweinen erkannt wurde. Die App zeigt
diese Information dann als verständliche Warnung an.

\folie{3}{Konzept}
Auf dieser Folie sieht man das Grundkonzept: Textiles Objekt,
Wahrnehmung und Reaktion.

Das textile Objekt ist der Teddybär. Er ist kindgerecht und vertraut.
Die Sensorik kann darin verdeckt eingebaut werden, sodass die Nutzung
keine große technische Hürde hat.

Die Wahrnehmung passiert über Geräusche. Der ESP32 misst den Pegel und
kann zusätzlich kurze Audiodaten an den Server schicken. Der Server
macht daraus ein Ereignis.

Die Reaktion passiert über die App. Statt nur rohe Messwerte zu zeigen,
übersetzt die App die Ereignisse in klare Zustände wie \glqq Alles
ruhig\grqq{}, \glqq Weinen erkannt\grqq{} oder \glqq Sehr laut\grqq{}.

\folie{4}{Architektur}
Die Architektur besteht aus drei Hauptteilen.

Der erste Teil ist der Teddybär mit ESP32 und Mikrofon. Der ESP32 misst
die Lautstärke und sendet die Daten per HTTP als JSON an den Server. Für
die Audioanalyse kann er außerdem kurze WAV-Ausschnitte senden.

Der zweite Teil ist der Go-Server. Er ist die zentrale Stelle im System:
Er nimmt Messwerte entgegen, speichert die letzten Ereignisse und
bewertet, ob eine Warnung ausgelöst werden soll. Die Bewertung kann
lokal über Lautstärke-Schwellwerte passieren oder optional mit Gemini
unterstützt werden.

Der dritte Teil ist die Flutter-App. Sie fragt den Server regelmäßig ab
und zeigt den aktuellen Status, den Pegel und die letzten Ereignisse an.
Dadurch müssen die Eltern nicht direkt auf den Bär schauen, sondern
bekommen die Warnung über die App.

\folie{5}{Flutter-App}
Die Flutter-App ist als Dashboard für Eltern gedacht.

Im oberen Bereich sieht man den aktuellen Status. Mögliche Zustände sind
zum Beispiel \glqq Bereit\grqq{}, \glqq Alles ruhig\grqq{}, \glqq
Weinen erkannt\grqq{} oder \glqq Sehr laut\grqq{}.

Zusätzlich zeigt die App den aktuellen Pegel an. Sie fragt den Server
alle drei Sekunden ab. Wenn ein neues Warnereignis erkannt wird, kann
ein Warnton abgespielt werden.

Unten gibt es eine Ereignisliste mit den letzten Meldungen. Für die Demo
sind außerdem Testbuttons hilfreich, weil man damit Warnungen auslösen
kann, auch wenn gerade kein ESP32 angeschlossen ist.

\folie{6}{Demo}
In der Demo würde ich zuerst kurz zeigen, dass der Server läuft. Danach
öffne ich die Flutter-App und zeige den normalen Zustand.

Dann löse ich ein Testereignis aus oder nutze den ESP32 mit Mikrofon.
Wenn ein lautes Geräusch oder ein Weinen-Ereignis erkannt wird, ändert
sich der Status in der App. Außerdem erscheint das Ereignis in der
Liste, und bei aktivierter Tonfunktion wird ein Warnsignal abgespielt.

Wichtig ist dabei: Der Prototyp funktioniert auch ohne externe KI. Wenn
Gemini nicht verfügbar ist, kann der Server weiterhin über lokale
Schwellwerte reagieren. Mit Gemini kann die Bewertung aber genauer
werden, weil nicht nur die Lautstärke, sondern auch die Art des Geräuschs
betrachtet wird.

\section*{Schluss}
Zusammenfassend zeigt das Projekt, wie ein normales textiles Objekt zu
einem interaktiven System erweitert werden kann.

Der Teddybär bleibt äußerlich ein vertrautes Objekt, bekommt aber durch
ESP32, Server und App eine neue Funktion. Er kann Geräusche wahrnehmen,
Ereignisse auslösen und Eltern über eine App informieren.

Für eine Weiterentwicklung könnte man die Hardware kleiner und stabiler
einbauen, die Erkennung mit mehr Testdaten verbessern und die App für
eine echte Nutzung auf dem Smartphone erweitern.

\section*{Kurzer Demo-Ablauf}
\begin{enumerate}[label=\arabic*.]
  \item Server starten.
  \item Flutter-App öffnen.
  \item Normalen Zustand \glqq Alles ruhig\grqq{} zeigen.
  \item Testbutton oder ESP32-Geräusch verwenden.
  \item Statuswechsel und Ereignisliste in der App zeigen.
  \item Kurz erklären, ob die Bewertung lokal oder mit Gemini erfolgt.
\end{enumerate}

\end{document}
